    En esta sección se describen los pasos a seguir para completar la tarea, la cual consiste en desarrollar código
    en C++ «\nameref{subsec:implementations}» y en realizar un informe. Se espera que cada
    uno de los pasos se realice de manera ordenada y siguiendo las instrucciones dadas.

    \begin{mdframed}
        Abra este documento en algún lector de PDF que permita hipervínculos, ya que en este documento el texto en
        color \textcolor{blue}{azul} suele indicar un hipervínculo.
    \end{mdframed}

\subsection{Implementaciones} \label{subsec:implementations}

\begin{enumerate}[(1)]
    \item Implementar cada uno de los algoritmos en C++:

    \begin{itemize}
        \item \textbf{Algoritmo Fuerza Bruta:} Resolver el problema mediante un enfoque exhaustivo que permita obtener una solución óptima para instancias pequeñas. \\
        Archivo: \texttt{code/implementation/algorithms/brute-force.cpp}.
        
        \item \textbf{Algoritmos Greedy:} Implementar dos heurísticas greedy (\textit{subóptimas}) para el problema. Estas heurísticas deben construir una selección válida de capítulos, aunque no necesariamente óptima. \\
        Archivos: \texttt{code/implementation/algorithms/greedy\{1,2\}.cpp}.

        Cada heurística greedy debe utilizar un criterio local distinto. Dicho criterio debe estar claramente descrito y justificado en el informe.
        
        \item \textbf{Programación Dinámica:} Implementar una solución óptima para el problema mediante programación dinámica. \\
        Archivo: \texttt{code/implementation/algorithms/dynamic-programming.cpp}.
    \end{itemize}

    \item \textbf{Restricciones de validez de una solución:}

    Una solución entregada por cualquier algoritmo debe cumplir las siguientes condiciones:

    \begin{itemize}
        \item Para cada anime, sólo se puede seleccionar un prefijo de capítulos.
        \item No se permite seleccionar un capítulo si no fueron seleccionados todos los capítulos anteriores del mismo anime.
        \item La suma total de minutos utilizados no puede exceder $M$.
        \item La suma total de energía utilizada no puede exceder $E$.
        \item El bono de completación de un anime se suma únicamente si se seleccionan todos sus capítulos.
    \end{itemize}

    \item \textbf{Mediciones de tiempo y memoria:} Implementar el programa principal en C++ \texttt{code/implementation/general.cpp} y su respectivo \texttt{makefile} que ejecute los algoritmos y genere archivos de salida en los directorios:

    \begin{itemize}
        \item \texttt{code/implementation/data/measurements/}
        \item \texttt{code/implementation/data/outputs/}
    \end{itemize}

    Los archivos generados deben permitir comparar el tiempo de ejecución, el uso de memoria y el resultado obtenido por cada algoritmo.

    \item \textbf{Generación de gráficos:} Implementar scripts en Python que lean los archivos de medición y generen gráficos en formato PNG en:

    \begin{itemize}
        \item \texttt{code/implementation/data/plots/}
        \item Scripts de ejemplo: \texttt{code/implementation/scripts/testcases\_generator.py} y \texttt{plot\_generator.py}.
    \end{itemize}

    Se espera generar gráficos que permitan analizar:

    \begin{itemize}
        \item Tiempo de ejecución según tamaño de entrada.
        \item Uso de memoria según tamaño de entrada.
        \item Calidad de solución de los algoritmos greedy respecto de la solución óptima.
        \item Comportamiento de los algoritmos al variar la cantidad de animes, capítulos, minutos disponibles y energía disponible.
    \end{itemize}

    \item \textbf{Documentación:} Documentar cada uno de los pasos anteriores:

    \begin{itemize}
        \item Completar el archivo \texttt{README.md} del directorio \texttt{code}.
        \item Documentar en cada uno de sus programas, al inicio de cada archivo, fuentes de información, referencias y bibliografía utilizada.
        \item Indicar claramente cómo compilar, ejecutar y generar los casos de prueba.
    \end{itemize}
\end{enumerate}

    Generar un informe en \LaTeX\ que contenga los resultados obtenidos y una discusión sobre ellos.

    \begin{itemize}
        \item No se debe modificar la estructura del informe.
        \item Las indicaciones se encuentran en el archivo \texttt{README.md} del repositorio y en la plantilla de \LaTeX.
        \item El informe final compilado (\texttt{reporte.pdf}) debe generarse dentro de la carpeta \texttt{report/}.
    \end{itemize}

    El informe debe incluir, al menos, los siguientes puntos:

    \begin{itemize}
        \item Explicación del algoritmo de fuerza bruta.
        \item Explicación de las dos heurísticas greedy implementadas.
        \item Explicación de la solución mediante programación dinámica.
        \item Análisis de complejidad temporal y espacial de cada enfoque.
        \item Comparación experimental de tiempo y memoria.
        \item Comparación de calidad de solución entre greedy y programación dinámica.
        \item Conclusiones sobre el comportamiento de los algoritmos.
    \end{itemize}